\chapter{Neural Architecture of Motor Control}
\label{ch:neural-arch}

\epigraph{%
  The motor system is not one controller.
  It is an orchestra of specialized modules,
  each playing its part in the symphony of movement.
}{}

\section{Cortical Motor Areas}

The cerebral cortex contains several distinct motor areas, each with specialized roles:

\subsection{Primary Motor Cortex (M1)}

M1 (Brodmann area 4) is the final cortical output stage for voluntary movement.
Layer V contains the giant pyramidal (Betz) cells whose axons form the corticospinal
tract---the most direct neural pathway from cortex to spinal motor neurons.

Key properties of M1:
\begin{itemize}
  \item \textbf{Somatotopic organization}: body parts are mapped to cortical regions
        (the motor homunculus), but with overlapping representations
  \item \textbf{Population coding}: individual neurons have broad tuning; movement
        direction is encoded by the \emph{population vector} \cite{Georgopoulos1986}
  \item \textbf{Preparatory activity}: M1 neurons are active 100--200~ms before
        movement onset, during motor preparation
\end{itemize}

\subsection{Premotor Cortex (PMC) and Supplementary Motor Area (SMA)}

\begin{itemize}
  \item \textbf{PMC}: movement planning based on external cues (visually guided reaching)
  \item \textbf{SMA}: movement planning based on internal cues (self-initiated sequences)
  \item Both areas project to M1 and directly to the spinal cord
\end{itemize}

\section{Subcortical Motor Structures}

\subsection{Cerebellum: The Forward Model}

The cerebellum contains more neurons than the rest of the brain combined ($\sim$70 billion
granule cells). Its primary role in motor control is as a \textbf{forward model}:

\begin{neurobox}{The Cerebellum as a Forward Model}{cerebellum-forward}
  The cerebellum:
  \begin{enumerate}
    \item Receives a copy of the motor command (efference copy) via the pontine nuclei
    \item Receives sensory feedback via the inferior olive and mossy fibers
    \item Computes the \emph{predicted} sensory consequences of the command
    \item Compares prediction with actual sensory feedback
    \item Generates an error signal that updates the forward model (climbing fiber
          pathway, long-term depression of parallel fiber--Purkinje cell synapses)
  \end{enumerate}
  In the language of control theory (Volume~I):
  \begin{itemize}
    \item The cerebellum implements $\hat{\state}_{k+1} = f(\hat{\state}_k, \control_k)$
          --- the state prediction step of the Kalman filter
    \item The climbing fiber error signal implements the innovation update
    \item The granule cell layer provides the basis functions for the nonlinear mapping
  \end{itemize}
\end{neurobox}

\subsection{Basal Ganglia: Action Selection}

The basal ganglia implement an action selection mechanism through competitive
inhibition. The direct pathway facilitates a desired action; the indirect pathway
suppresses competing actions. This is analogous to a model predictive controller
evaluating multiple candidate trajectories and selecting the best one.

\section{Spinal Cord: The Local Controller}

The spinal cord is not merely a relay station. It contains:
\begin{itemize}
  \item \textbf{Motor neuron pools}: lower motor neurons that directly innervate muscles
  \item \textbf{Interneuron networks}: local circuits that implement reflexes, pattern
        generation, and coordination
  \item \textbf{Central pattern generators} (CPGs): oscillatory circuits for rhythmic
        movements (Chapter~8)
\end{itemize}

\section*{Exercises}

\begin{enumerate}
  \item \textbf{Population Vector.}
        Given neurons with preferred directions every $45^\circ$, compute the population
        vector for a movement at $30^\circ$. Plot individual neuron contributions.

  \item \textbf{Cerebellar Forward Model.}
        Implement a simple forward model using a lookup table. Show how climbing
        fiber errors update the model over successive trials.

  \item \textbf{(Research.)} Read Georgopoulos et al.\ (1986) and summarize the
        evidence for population coding in M1.
\end{enumerate}
